\chapter{Results}\label{ch:results}

This chapter reports what was measured, in mechanism order: first what
the judge does, then what the partition does. The order matters
because the strongest result in this project is a measurement of the
judge, and it is what makes the partition results of
Section~\ref{sec:res-partition} interpretable. The judge bypasses the
partition on this corpus by verifying answers against a key it was
never shown, and the cell-scoped rubric does not change the ranking
that judging produces: both arms recover the same ordering in every
domain, and what separates them is how often the cell-scoped arm
declines to name a winner. A reader who learns in
Section~\ref{sec:res-judge} what the judge conditions on already knows
why that is a property of the corpus and of the way this judge works,
not a failure of the construction.

The reading conventions below govern every result in this chapter;
each is stated once here and not repeated at the point of use.

\begin{center}
\fbox{\parbox{0.92\textwidth}{%
\textbf{Reading conventions for every result in this chapter.}
\begin{itemize}
\item \textbf{Tie policy.} Every Spearman $\rho$ appears under both
tie conventions --- ties half-weighted, then tied comparisons dropped,
in that order --- per rule~D5 of Section~\ref{sec:meth-discipline};
the measured instability that forces the rule is
Section~\ref{sec:res-tie-policy}. Every table keeps both rows.
\item \textbf{Estimator.} Bradley--Terry $\theta$ is quoted, never
win-rate, because the scheduler over-samples uncertain pairs and
win-rate is biased by opponent strength
(Section~\ref{sec:meth-scheduling}). The over-sampling is measurable.
Across the settled batch the most-scheduled model appears in
$28.6\%$ to $30.9\%$ of a run's comparisons, against the $25.0\%$ a
uniform scheduler would give at eight models. A win-rate computed on
that allocation is not comparable across models.
\item \textbf{Seed.} All construction and arena results are at seed
42. There is no three-seed convention behind any headline number in
this thesis, and none is claimed.
\item \textbf{Counts.} Comparison counts are logical,
$B_{\text{logical}}$; each logical comparison issues two judge calls
under the dual-call design, so $B_{\text{API}} = 2B_{\text{logical}}$.
\item \textbf{Quantisation.} Spearman $\rho$ over $M$ models moves on
a grid of $6/(M(M^2-1))$: $0.0119$ at $M = 8$, $0.0046$ at $M = 11$,
$0.0035$ at $M = 12$. Differences below one grid step are not
differences.
\item \textbf{Sign of a paired contrast.} Throughout this chapter
$\Delta\rho$ is \textsc{main} $-$ \textsc{c5}, positive favouring
the cell-scoped arm. It is reported for continuity only
(Section~\ref{sec:exp-metrics}): its seed-to-seed spread exceeds the
withdrawn decisive threshold, and no decision rests on it.
\end{itemize}}}
\end{center}

Interpretation is deferred to Chapter~\ref{ch:discussion}.

Because this thesis withdrew two run generations and corrected its own
apparatus several times, the status of every evidence source is fixed
once, here, and not re-litigated at each point of use
(Table~\ref{tab:evidence-status}). A number in this chapter is settled
unless this table says otherwise.

\begin{table}[htbp]
\centering
\small
\caption{Status of every evidence source. ``Settled'' numbers carry no
inline caveat when quoted; the others are quoted only with their
status named.}
\label{tab:evidence-status}
\begin{tabular}{p{5.3cm}p{2.3cm}p{5.2cm}}
\toprule
Evidence source & Status & Where governed \\
\midrule
Settled batch (8 models, 8 domains, \texttt{RESOLVED}) &
\textbf{settled} & Table~\ref{tab:run-record} \\
Construction artifact + fixed-point certificate &
settled & \S\ref{sec:res-construction} \\
Offline nulls: granularity, leaf-substructure, rubric-specificity,
$J$-floor, screen refit & settled (GT-only / offline) &
\S\ref{sec:res-partition}; rosters named per table \\
Pilot arena (Run A) & \textbf{withdrawn} & \S\ref{sec:res-capture-gap};
survives only as the defect record \\
Twelve-model paired campaign & \textbf{withdrawn} &
\S\ref{sec:res-runs}; survives only in registrations
(Appendix~\ref{app:preregs}) \\
Registered $|\Delta\rho| \geq 0.007$ decision rule & \textbf{withdrawn}
& Appendix~\ref{app:preregs}, addendum to P4 \\
Bradley--Terry confidence intervals & \textbf{none reported} &
variance repair postdates the stored errors
(Section~\ref{sec:disc-debt}) \\
Reliability constant $c = 2.52$ & unrefitted (11-model) &
Section~\ref{sec:disc-debt} \\
Routing ECE export ($0.2114$) & diagnostic only &
Appendix~\ref{app:metrics} \\
\bottomrule
\end{tabular}
\end{table}

\section{What Was Run}\label{sec:res-runs}

One body of arena evidence and seven offline analyses produced the
results in this chapter: the settled eight-domain batch, to which every
arena number in this thesis is attributed. Two earlier bodies of runs
are absent. A pilot Math arena ran on a corpus whose stored responses
were incomplete --- for four of its eight models no response text
survived, and format rather than content decided roughly half of its
comparisons --- and a twelve-model paired campaign ran before two
scheduler repairs that the defect analysis of the pilot and its
successors forced. Both are withdrawn in their entirety
(Table~\ref{tab:evidence-status}); they survive only in the
registration record of Appendix~\ref{app:preregs} and the corrections
register.

\textbf{The settled batch --- eight paired domains on the eight-model
roster.} Snapshot \texttt{20260727\_042523}, the frozen artifact whose
construction evidence is Section~\ref{sec:res-construction}. Eight
models (Table~\ref{tab:roster-r8},
Appendix~\ref{app:models}), selected to span ground-truth accuracy
$16.5\%$ to $90.5\%$ so that most model pairs are key-decidable at the
question counts a domain routes. Judged in \texttt{RESOLVED} option
mode: the judge receives the question stem without its options block,
and each response carries its model's selected option attached as text
(attachment $97.6\%$ pooled; the per-domain range is stated once in
Chapter~\ref{ch:experimental-design}). Two conditions per domain:
\textsc{main}, each question judged under its own cell's rubric, and
C5 (\textsc{generic}), which replays \textsc{main}'s exact
$(\text{question}, \text{model pair})$ triples under the same cells and
a generic MT-Bench-style pairwise prompt. One run per domain,
sequentially, on 2026-07-31; no run crashed, no verdict was removed by
the confidence gate, and no rate-limit response occurred in any domain.
Table~\ref{tab:run-record} is the run record; the results derived from
it are Section~\ref{sec:res-c5} and the sections that cite them.

\begin{table}[htbp]
\centering
\caption{Run record of the settled batch. Verdicts are per arm;
\textsc{generic} replays \textsc{main}'s triples exactly, short by
three comparisons on law and one on engineering, where a failed replay
was not retried.}
\label{tab:run-record}
\begin{tabular}{lrrr}
\toprule
Domain & Routed questions & Cells & Verdicts per arm \\
\midrule
mathematics      & 291 & 11 & 1{,}030 \\
physics          & 229 & 10 & 455 \\
law              & 276 & 6  & 811 \\
engineering      & 158 & 4  & 283 \\
psychology       & 265 & 8  & 1{,}130 \\
philosophy       & 129 & 3  & 431 \\
history          & 131 & 4  & 495 \\
computer science & 158 & 6  & 526 \\
\midrule
total            & 1{,}637 & 52 & 5{,}161 \\
\bottomrule
\end{tabular}
\end{table}

\begin{figure}[htbp]
\centering
\includegraphics[width=\textwidth]{rank_recovery_grid}
\caption{The batch on one page: each model's arena rank displacement
from its key rank, per domain. White is exact agreement ($45$ of $64$
cells); upward displacement is blue (deeper blue for the single $+2$,
psychology), downward dark grey, and every cell prints its signed
offset, so the encoding survives grayscale. All displacements come
from two recurring near-tie pairs (Llama-3.1-70B against Qwen1.5-72B; arx\_0314 against
gpt-4o). A small circle marks a model that sits in at least one
key-undecidable pair in that domain --- there the key itself does not
fully constrain its position. Each column reproduces its domain's
published $\rho$ exactly (asserted by
\texttt{make\_rank\_recovery\_grid.py}).}
\label{fig:rank-recovery-grid}
\end{figure}

\textbf{What ``paired'' means here, exactly.} The second arm replays
the first arm's comparison list --- the same question, the same model
pair, the same cell, for every comparison the first arm judged. The two
arms therefore share questions, comparisons and cells by construction,
not statistically: measured as the overlap of
$(\text{question}, \text{model pair})$ tuples, pairing is $100\%$ in
six domains and $99.6\%$ on law and engineering. What the arms do not
share is the judge prompt and the verdict format, and
Table~\ref{tab:arm-differences} (Appendix~\ref{app:judge-generation})
itemises exactly what travels with that difference.

\textbf{The offline analyses} required no judge calls: the fixed-point
certificate and the acceptance-gate seed sweep
(Section~\ref{sec:res-construction}); the rubric-specificity null
(Section~\ref{sec:res-rubric-null}); the pre-registered granularity
analysis (Section~\ref{sec:res-granularity}); the leaf-substructure
permutation null (Section~\ref{sec:res-leaf-null}); the
domain-reordering screen (Section~\ref{sec:res-law}); the cross-link
proposal analysis that retired polyhierarchy
(Section~\ref{sec:poly-negative}); and the reliability-constant refit
reported with the measurement debt in Section~\ref{sec:disc-debt}.

\section{What the Judge Does}\label{sec:res-judge}

\subsection{The Judge Decides on Answer Correctness}\label{sec:res-correctness}

A model pair is \emph{key-decidable} in a domain at
$N$ routed questions when the two accuracies separate by more than
twice their pooled standard error,
$|p_1 - p_2| > 2\sqrt{p_1(1-p_1)/N + p_2(1-p_2)/N}$. On the settled
batch the key can decide $200$ of the $8 \times 28 = 224$ domain-level
pairs, and the arena's Bradley--Terry ordering agrees with the key on
\textbf{196 of those 200} ($98.0\%$) --- one configuration, no
per-domain tuning (Table~\ref{tab:settled-recovery}). All four misses
sit at the smallest accuracy gaps in their domains, where the key
itself is nearly indifferent.

\begin{table}[htbp]
\centering
\caption{Rank recovery and judge--key agreement on the settled batch,
\textsc{main} arm. $\rho$ is Spearman against the ground-truth accuracy
ordering and is \emph{identical under both tie conventions} in every
domain and both arms, so one column carries it. Agreement is on
key-decidable comparisons where the judge committed to a winner.}
\label{tab:settled-recovery}
\begin{tabular}{lrrrr}
\toprule
Domain & $\rho$ & decidable pairs & recovered & agreement \\
\midrule
mathematics      & 1.0000 & 28/28 & 28/28 & 89.5\% \\
physics          & 0.9762 & 26/28 & 25/26 & 89.0\% \\
law              & 0.9524 & 25/28 & 23/25 & 78.9\% \\
engineering      & 1.0000 & 26/28 & 26/26 & 86.6\% \\
psychology       & 0.9286 & 24/28 & 23/24 & 87.2\% \\
philosophy       & 0.9762 & 24/28 & 24/24 & 90.3\% \\
history          & 0.9524 & 22/28 & 22/22 & 80.4\% \\
computer science & 0.9524 & 25/28 & 25/25 & 86.6\% \\
\midrule
pooled           &        & 200/224 & \textbf{196/200} & 86.0\% \\
\bottomrule
\end{tabular}
\end{table}

\begin{figure}[htbp]
\centering
\includegraphics[width=\textwidth]{settled_validity_scatter}
\caption{The arena tracks the key wherever the key can see. One point
per (domain, model pair), $n = 224$: the pair's ground-truth accuracy
gap against its Bradley--Terry score difference, signed so that
agreement with the key is positive. Filled points are the $200$
key-decidable pairs; small grey diamonds the $24$ undecidable ones.
The grey band extends to the largest gap that is undecidable in any
domain (the criterion is per-domain, so the band is a summary, not
the test). The four misses --- in blue, each labelled --- all sit at
the smallest decidable gaps, where the key itself is nearly
indifferent.
Generated by \texttt{make\_settled\_validity\_scatter.py}, which
asserts the $224/200/196/4$ counts before drawing.}
\label{fig:settled-validity}
\end{figure}

\textbf{What 196/200 is, and what it is not.} The ground-truth ordering
the eight domains are scored against is \emph{identical} in seven of
them; only law differs, by one adjacent swap of two nearly-tied models
($0.207$ against $0.217$). The roster was selected for a large global
accuracy spread, and that spread survives any subsample of the corpus.
The batch therefore demonstrates that the arena recovers \emph{one}
ordering, eight times, on disjoint question samples with independent
verdicts --- a real replication, but of one result, not eight. Every
use of the $196/200$ figure in this thesis carries this qualification.

Where the judge commits to a winner on a
key-decidable comparison, it agrees with the key on $86.0\%$ pooled
($78.9\%$ to $90.3\%$ by domain). Two qualifications travel with these
numbers. First, agreement is reported over the decidable subset only
--- on an undecidable comparison there is no accuracy-derived winner to
agree with --- and conditional on the judge committing; the abstention
behaviour is Section~\ref{sec:res-tie-signature} and the
committed-versus-abstained decomposition is Section~\ref{sec:res-c5}.
Second, the judge is answer-key-blind but not answer-blind: in the
settled runs the options block is withheld, but $97.6\%$ of stored
traces state the answering model's own selection
(Section~\ref{sec:arch-judge}). Nothing in the prompt tells the judge
which option is right; a judge able to solve the item does not need to
be told.

\subsection{The Tie Rate Follows the Key}\label{sec:res-tie-signature}

Not part of the hypothesis: the judge's tie
rate nearly doubles exactly where the answer key stops discriminating
(Table~\ref{tab:tie-by-key}). On comparisons the key can decide, the
\textsc{main} arm ties on $13.0\%$; on comparisons where both responses
are right or both are wrong, $24.0\%$. The direction holds in all eight
domains (ratios $1.3$ to $2.5$), and the generic arm shows the same
signature from a lower base ($10.9\%$ decidable, $15.1$--$27.7\%$
undecidable by domain).

\begin{table}[htbp]
\centering
\caption{\textsc{main}-arm tie rate by whether the answer key can
discriminate between the two responses, settled batch.}
\label{tab:tie-by-key}
\begin{tabular}{lrrr}
\toprule
Domain & key-decidable & key-undecidable & ratio \\
\midrule
mathematics      & 11.4\% & 24.1\% & 2.1 \\
physics          &  9.9\% & 24.3\% & 2.5 \\
law              & 15.3\% & 26.5\% & 1.7 \\
engineering      & 14.7\% & 22.6\% & 1.5 \\
psychology       & 12.5\% & 25.2\% & 2.0 \\
philosophy       & 14.8\% & 27.7\% & 1.9 \\
history          & 12.6\% & 18.3\% & 1.5 \\
computer science & 14.5\% & 19.1\% & 1.3 \\
\midrule
pooled           & \textbf{13.0\%} & \textbf{24.0\%} & 1.8 \\
\bottomrule
\end{tabular}
\end{table}

A near-doubling of abstention precisely where the key stops helping is
the signature of a correctness-driven decision: when the judge can
verify a difference it names a winner, and when there is no difference
to verify it declares equivalence far more often. The undecidable
stratum also decomposes asymmetrically --- $28.7\%$ ties where both
responses are correct against $20.8\%$ where both are wrong --- the
direction \textcite{stephan2024calculation} predict, since two
verifiably right answers are easier to call equivalent than two
wrong ones. The tie rate is
recorded for the Bradley--Terry sufficient statistics, not as a test of
anything, which is why this counts as corroboration rather than as
fit. Figure~\ref{fig:tie-signature} shows the same sixteen rates as
Table~\ref{tab:tie-by-key} and the tie columns of
Table~\ref{tab:c5-decisiveness} in one view: eight dumbbells, all
leaning the same way, in both arms.

\begin{figure}[htbp]
\centering
\includegraphics[width=\textwidth]{tie_signature}
\caption{The tie signature, per domain and pooled: the judge's tie
rate on key-decidable comparisons (circle) against key-undecidable
ones (arrowhead), cell-scoped arm filled, generic arm hollow,
right-margin labels giving the cell-scoped ratio. The direction holds
in every domain and both arms; the generic arm shows the same
signature from a lower base. Generated by
\texttt{make\_tie\_signature.py}, which asserts the pooled rates
before drawing.}
\label{fig:tie-signature}
\end{figure}


\subsection{Verification, and Its Rival: A Competence Prior}\label{sec:res-competence}

The strongest alternative reading of everything above comes from the
literature this thesis extends: \textcite{stephan2024calculation} show
that judges favour the generally stronger model even when it is wrong,
and note that on aggregate ranking recovery this bias and genuine
verification predict the \emph{same} outcome. The $196/200$ of
Section~\ref{sec:res-correctness} therefore cannot, alone, distinguish
the two. Three offline measurements on the settled batch's stored
verdicts discriminate.

First, a judge that \emph{always} picks the pair-stronger model scores
$83.5\%$ on the $2{,}978$ key-decidable comparisons, against the real
judge's $86.0\%$ committed agreement: the aggregate headroom over pure
reputation is small.

Second, and decisively, splitting decidable comparisons by which side
holds the correct answer: where the correct answer sits with the
pair-\emph{stronger} model the judge agrees with the key on $91.1\%$
of committed verdicts ($n = 2{,}208$); where it sits with the
pair-\emph{weaker} model, $56.5\%$ ($n = 382$). The judge is far above
the near-total failure \textcite{stephan2024calculation} report in the
same split (often below $20\%$), so item-level verification is real
--- but the $34$-point collapse shows a competence prior sharing the
verdict with it.

Third, the tie signature of Section~\ref{sec:res-tie-signature} is the
behaviour a prior cannot produce: a reputation heuristic has no reason
to double its abstention exactly where the key stops discriminating,
because the responses' lengths, styles and authors are the same on
both sides of that boundary.

Corrected claim this section leaves standing:
the verdicts are produced by \textbf{verification and a competence
prior jointly}. Verification is visible wherever the two decouple; the
prior is visible in the same place. Every aggregate number in this
chapter is co-produced by both, and the roster's wide spread ---
chosen to make pairs decidable --- is also what keeps the two aligned
on $2{,}208$ of $2{,}590$ committed decidable comparisons. For
cross-paper comparison, the unconditional decidable agreement
(abstention scored as a miss) is $74.8\%$, inside the $57$--$80\%$
band \textcite{stephan2024calculation} report for their strongest
judges on the analogous subset.

\subsection{The Capture Gap: Why the Pilot Is Withdrawn}\label{sec:res-capture-gap}

This subsection reports a defect and its repair, not a result. It is
what withdrew the pilot run, and it is kept because five other parts of
this thesis --- the exclusion policy, the corpus repair, the roster
constraints, and two scope statements --- rest on it.

Four of the pilot's eight models had their stored
response text recorded as the empty string across all 12{,}032 rows,
while the predicted answer field was populated: a capture gap in the
evaluation-generation pipeline, not model behaviour. The judge was
therefore shown a worked solution on one side and a synthesised
one-line stub on the other, and on such mixed pairs it took the side
with text on $863$ of $868$ decided comparisons --- half of all the
pilot's comparisons. Those verdicts measure the presence of text, not
model quality, and every number derived from that run is withdrawn
rather than qualified.

Repair: the cause was in the ingest loader,
which had dropped stored response text for 34 of the corpus's 47
evaluated models. Re-ingest backfilled 356{,}402 rows and recovered 24
of the 34, taking the corpus from 13 models with usable traces to 37;
ten had no persisted text to recover, and one of the 37 sits in a
different question-identifier space and is excluded
(Appendix~\ref{app:models}), so the maximum valid roster is 36 models.
The settled batch runs entirely on the repaired corpus, and its roster
admits no model whose stored output is majority-bare --- though one,
at $21.4\%$ bare, sits under that policy's threshold rather than clear
of it (Appendix~\ref{app:models}).

\subsection{Off-Subject Questions Are Judged Like On-Subject Ones}\label{sec:res-offsubject}

Routing is subject-coherent but not perfectly so:
$77.8\%$ of the $1{,}637$ judged questions carry the MMLU-Pro category
their cell was built from ($68.4\%$ engineering to $94.6\%$ law), and
part of the $22.2\%$ residual is the benchmark's own label taxonomy
rather than misrouting --- graph-theory combinatorics filed under
computer science lands, correctly, in a combinatorics cell. The
remainder is genuinely off-subject material judged under another
subject's rubric: a chemistry item under a combinatorics rubric is the
concrete case.

\textbf{the judge does not care.} Within cells that received at least ten key-decidable
comparisons on each side ($27$ cells qualify), agreement with the key
on off-subject questions differs from on-subject ones by a median of
$+0.95$ percentage points, with $14$ cells favouring on-subject and
$13$ favouring off-subject (sign test $p = 1.0$); the tie rate is
likewise indistinguishable ($15$ against $12$, $p = 0.70$). A rubric
written for extremal combinatorics neither helps nor hinders the judge
on the alkene question that strayed into its cell.

Two boundaries on this result. It is a null with modest power --- only
$27$ cells clear the ten-per-side floor, so ``no detectable effect''
is the claim, not ``no effect''. And it varies the \emph{material}
while Section~\ref{sec:res-c5} varies the \emph{prompt}; the two
measurements agree that on this corpus the cell rubric does not decide
verdicts, and each covers the other's blind side.

The continuous version of the same question replaces the
on/off-subject dichotomy with distance
(Figure~\ref{fig:rubric-reach}). Measured as cosine distance to the
centroid of the queries its rubric was induced from, a routed held-out
question sits at $0.511$ --- $38\%$ of the way from a random reserved
question ($0.607$) toward construction-grade proximity ($0.353$), and
only $6$--$21\%$ of the way in five of eight domains. The premise of
the rubric-coverage worry is therefore real: rubrics are applied well
outside the material they were induced from. The consequence is not:
within cells, one standard deviation of that distance moves the
cell-scoped arm's tie rate by $+1.02$ percentage points --- and moves
the \emph{generic} arm, which carries no domain criteria at all, by
$+1.06$. The within-triple difference is $-0.04$ points
($95\%$ CI $[-1.39, +1.30]$, cell fixed effects, $n = 2{,}788$).

\begin{figure}[htbp]
\centering
\includegraphics[width=\textwidth]{rubric_reach}
\caption{The rubric is applied far outside its induction pool, and
the judge does not care. Panel~(a): per domain, the cosine distance
from a routed held-out question to its rubric's induction-pool
centroid (dot), between construction-grade proximity (left rule) and
a random held-out question (right rule). Panel~(b): change in tie
rate and key agreement per standard deviation of that distance, with
cell fixed effects; the arm with no rubric moves identically, so the
effect is not rubric coverage. Intervals are regression confidence
intervals, not Bradley--Terry intervals. Generated by
\texttt{make\_rubric\_reach.py}.}
\label{fig:rubric-reach}
\end{figure}

\subsection{Tie Policy, Sampler, and the Stability of \texorpdfstring{$\rho$}{rho}}\label{sec:res-tie-policy}

It bounds every rank correlation in this thesis.
On the settled batch the tie convention moves nothing: half-weighting
ties and dropping tied comparisons produce the identical $\rho$ in
every domain and both arms (Table~\ref{tab:settled-recovery}), because
the roster's accuracy spread leaves no adjacent pair close enough for
the convention to flip. The scheduler seed is the live source of
movement: two runs of an identical configuration --- same domain,
taxonomy, question pool, option mode, differing only in the scheduler
seed --- moved the paired contrast $\Delta\rho$ by six and zero grid
steps ($0.0714$ and $0.0000$ at $M = 8$), while per-verdict answer-key
agreement reproduced to within $0.1$ percentage points across the same
two runs.

Absolute $\rho$ from this pipeline is not stable at
the precision a three-decimal figure implies, and no decision in this
thesis rests on a $\rho$ difference; the per-verdict statistics carry
the power. This is the evidence behind rule~D5 of
Section~\ref{sec:meth-discipline} and behind the withdrawal of the
registered $\Delta\rho$ threshold
(Appendix~\ref{app:preregs}).

\subsection{Rubric Specificity at Matched Budget}\label{sec:res-c5}


This section compares two judge prompts under the replay design of
Section~\ref{sec:res-runs}: the cell-scoped rubric induced for each
cell against a generic MT-Bench-style pairwise prompt on the identical
comparison set. Neither arm varies the partition
(Appendix~\ref{app:preregs}, addendum to P4), and the two arms differ
in more than the rubric text: system prompt, user scaffold, verdict
format and confidence handling all travel together
(Table~\ref{tab:arm-differences},
Appendix~\ref{app:judge-generation}). What is measured is the
cell-scoped judging condition as a bundle, at matched budget.

Rankings: the contrast is empty. Both arms produce
the same Bradley--Terry ordering in seven of eight domains, under
both tie conventions; in computer science they differ by one adjacent
transposition ($\Delta\rho = 0.0238$, two grid steps --- inside the
six-step seed-to-seed spread of
Section~\ref{sec:res-tie-policy}, and not read as an effect). A
rubric written for the cell does not change which model this judge
ranks higher.

Verdicts: what separates the arms is decisiveness
(Table~\ref{tab:c5-decisiveness}). On key-decidable comparisons the
cell-scoped arm abstains on $13.0\%$ against the generic arm's
$10.9\%$, and the direction is the same in all eight domains.
Conditional on committing to a winner, the two agree with the key at
$86.0\%$ and $87.2\%$ pooled --- about one point apart. On the $37$
comparisons where both arms committed and exactly one matched the
key, the split is $11$ against $26$ in the generic arm's favour
(two-sided binomial $p = 0.020$), significant pooled and in no single
domain.

\begin{table}[htbp]
\centering
\caption{The two arms on shared key-decidable comparisons, per
domain: abstention rate, agreement with the key conditional on
committing, and the discordant comparisons where both arms committed
and exactly one was right.}
\label{tab:c5-decisiveness}
\begin{tabular}{lrrrrr}
\toprule
& \multicolumn{2}{c}{tie rate} & \multicolumn{2}{c}{agreement $\mid$ committed} & \\
\cmidrule(lr){2-3}\cmidrule(lr){4-5}
Domain & \textsc{main} & \textsc{gen} & \textsc{main} & \textsc{gen} & discordant \\
\midrule
mathematics      & 11.4\% &  9.5\% & 89.5\% & 91.4\% & 1/5 \\
physics          &  9.9\% &  7.9\% & 89.0\% & 90.6\% & 1/2 \\
law              & 15.3\% & 14.1\% & 78.9\% & 81.2\% & 2/7 \\
engineering      & 14.7\% & 10.7\% & 86.6\% & 85.7\% & 0/0 \\
psychology       & 12.5\% & 11.2\% & 87.2\% & 88.5\% & 2/3 \\
philosophy       & 14.8\% & 11.5\% & 90.3\% & 86.5\% & 3/2 \\
history          & 12.6\% &  9.4\% & 80.4\% & 81.8\% & 2/4 \\
computer science & 14.5\% & 11.7\% & 86.6\% & 89.4\% & 0/3 \\
\midrule
pooled           & \textbf{13.0\%} & \textbf{10.9\%} & 86.0\% & 87.2\% & \textbf{11/26} \\
\bottomrule
\end{tabular}
\end{table}

A rubric written for a cell changes how often this
judge abstains, not which model it ranks higher. The abstention
difference is attributable to the condition as a whole, not to the
rubric text alone, because the uncontrolled verdict-format difference
(a structured tie field against a trailing tie token) points at
exactly the quantity that moved. The reasoning scaffold runs in the
opposite direction (the cell-scoped arm is the one forced through a
four-step critique before its verdict, and it is the \emph{less}
decisive arm), so scaffolding is an unlikely explanation; format and
confidence handling remain uncontrolled. The isolation run that would
attribute the effect --- the generic system prompt through the
cell-scoped arm's own template and parser, replaying these same
triples --- is one condition and one domain, and is named as future
work (Section~\ref{sec:disc-future}).

For completeness, under the convention that scores an abstention as a
miss, pooled decidable-comparison agreement is $74.8\%$ against
$77.7\%$; that blend of decisiveness and accuracy is reported only
here, and no claim rests on it.
Figure~\ref{fig:discordance-split} separates the two effects the
convention would blend.

\begin{figure}[htbp]
\centering
\includegraphics[width=\textwidth]{discordance_split}
\caption{RQ2's two effects, kept apart. Panel~(a): tie rate on
decidable comparisons per domain, cell-scoped (filled) to generic
(hollow) --- eight slopes, one direction; this is the robust effect.
Panel~(b): the $37$ both-committed discordant comparisons split
$11$ against $26$, with the exact $95\%$ interval of the generic-only
count clearing the even split ($p = 0.020$) --- real, thin, and
significant in no single domain. Generated by
\texttt{make\_discordance\_split.py}.}
\label{fig:discordance-split}
\end{figure}


\subsection{Identification Cost}\label{sec:res-cost}

On the settled batch, from each run's comparison
log in arrival order. Per-cell \emph{identification} --- a connected
comparison graph over the eight-model roster --- is reached in
\textbf{all 52 cells}, at between $7$ and $10$ comparisons per cell,
median $7$. Seven is the graph-theoretic floor (a spanning tree over
eight models), so the scheduler identifies a cell at or within three
comparisons of the cheapest possible schedule. The measured cost sits
far inside the $\leq 40$-comparison band fixed in advance
(Appendix~\ref{app:preregs}, P3). The registration's void condition
--- a cost \emph{constant} across cells would indicate the scheduler
was still capping --- is not triggered by a cost constant at the
structural floor: a cap mimics a constant above the floor, not the
floor itself.

This result pairs with its opposite
(Figure~\ref{fig:budget-vs-resolution}). Identification is nearly free
--- $52$ cells for roughly $400$ logical comparisons --- but
\emph{resolution} is not: separating a cell's adjacent model pairs at
the corrected standard errors requires on the order of $143$
comparisons per pair, and the batch supplies a median of $1.6$ to
$5.1$ (Section~\ref{sec:res-granularity}): two orders of magnitude
short, at this budget, in every cell of every domain.

\begin{figure}[htbp]
\centering
\includegraphics[width=\textwidth]{budget_vs_resolution}
\caption{Identification is free and resolution is not. One dot per
cell: comparisons per model pair achieved in that cell, log scale.
The grey band is the identification cost (the whole eight-model graph
connects within $7$--$10$ comparisons, i.e.\ $0.25$--$0.36$ per
pair); the dashed rule is the $\sim 143$ comparisons per pair needed
to resolve an adjacent pair. Every cell in every domain sits in the
gulf between the two. No intervals are drawn, per the no-interval
rule of Section~\ref{sec:exp-metrics}. Generated by
\texttt{make\_budget\_vs\_resolution.py}.}
\label{fig:budget-vs-resolution}
\end{figure}

\subsection{What Section~\ref{sec:res-judge} Establishes}\label{sec:res-judge-summary}

All on the settled batch. The judge
\textbf{decides on answer correctness}: it recovers the key's ordering
on $196$ of $200$ decidable pairs and agrees with the key on $86.0\%$
of decidable comparisons where it commits
(Section~\ref{sec:res-correctness}); its \textbf{tie rate nearly
doubles} exactly where the key stops discriminating, in all eight
domains (Section~\ref{sec:res-tie-signature}); its \textbf{agreement
splits on candidate strength} --- $91.1\%$ where the correct model is
the pair's stronger, $56.5\%$ where it is the weaker
(Section~\ref{sec:res-competence}); it scores
\textbf{off-subject questions like on-subject ones}, so the cell
rubric's domain criteria are not what decides
(Section~\ref{sec:res-offsubject}); and replacing the cell rubric
with a generic prompt \textbf{changes its abstention rate, not its
ranking} (Section~\ref{sec:res-c5}). Five measurements; the last two
are matched controls that vary the material and the prompt
respectively while sharing every other confound.

All of the above: \textbf{on a multiple-choice corpus with a
verifiable key, a capable judge's verdicts are produced jointly by
answer verification and a competence prior --- it shortcuts to
correctness where it can check, leans on candidate strength where it
cannot, and by either route bypasses the partition entirely.}

One alternative explanation is dismissed explicitly, because the
settled roster invites it. The roster is not length-matched --- median stored response spans
$6.79\times$, and across the 36 valid models length correlates with
accuracy at $r = +0.597$ (Appendix~\ref{app:models}) --- so a judge
that merely preferred longer responses would also recover much of the
key's ordering on this roster. The tie signature is what separates the
two accounts. A length-preferring judge has no reason to double its
abstention rate exactly where the answer key stops discriminating: the
responses are the same lengths on both sides of that boundary, which
is drawn by the key, not by the prompt. A correctness-verifying judge
doubles it necessarily, because where both answers check out (or
neither does) there is nothing left to verify. Table~\ref{tab:tie-by-key}
therefore rules out length preference as the primary mechanism; what
the length confound still contaminates is the \emph{margin} of
individual verdicts, and that is carried as a limitation
(Section~\ref{sec:disc-debt}).

This is a hypothesis with strong support, not a demonstration, and
its evidence sits on the middle rung of a three-rung ladder. Under
\texttt{OPTIONS} (the withdrawn pilot's mode) the judge saw the full
options block. Under \texttt{RESOLVED} (every settled run) the
options block is withheld but $97.6\%$ of responses state their own
selection (Section~\ref{sec:res-correctness}). The registered
options-blind arm (P5) --- no options \emph{and} responses stripped of
their stated answers --- has not been run; its predictions are fixed
in writing (Appendix~\ref{app:tuning-protocol}) and it is blocked on
a trace-complete roster, because without options and without stored
reasoning an untraced-versus-untraced comparison is uninterpretable
rather than a chance-level null. The settled batch narrows the
verification channel; it does not close it.

\section{What the Partition Does}\label{sec:res-partition}

\subsection{Link 1: The Construction Meets Its Requirements}\label{sec:res-construction}

Frozen artifact has 154 nodes, 87 leaves,
maximum depth 6, $J = 0.253129$, and conserved query mass $8299.00$
with no mass warnings. All 87 leaves are labelled and carry a judge
rubric. Leaf populations have median 88 queries, IQR $[71, 118]$,
minimum 54 and maximum 396 (Figure~\ref{fig:leaf-size-hist}), with one
leaf ending below the birth floor of 55 for the reason given in
Section~\ref{sec:arch-limits}.

Construction reaches a certified fixed point at
iteration 10 on all four terms, every one at or below $10^{-12}$
against a tolerance of $10^{-6}$ (the exact certificate values are in
Appendix~\ref{app:dag-formalism}); iterations 8, 9 and 10 are
identical on every recorded column. The re-routing term is
the one that cannot be dropped: it is the only operation in the loop that
moves the objective without passing through the acceptance gate. One further
application of the operator would change nothing the artifact consists
of: not the structure, not the partition, and not the fitted node
parameters, which matters because held-out queries route through those
parameters. The certificate tests period~1 only: it asks what one more
iteration would do, not what any longer cycle would do.

It corrects how the stop is usually described.
The loop's stopping rule is a disjunction of two independent criteria
(Section~\ref{sec:arch-limits}). The structural arm requires five
consecutive iterations with zero graph edit distance \emph{and} a
minimum iteration count of $\max(0.8 \times \text{anchors}, 5)$, which
is $11$ at fourteen anchors, so it could not have fired at iteration
$10$. The objective arm requires five consecutive iterations in which
both the edit-phase and the re-routing changes in $J$ sit within
$\tau = 10^{-6}$. The frozen run's log closes with ``Early stopping
triggered in iteration 10 due to convergence (J stationarity)'' and
reports graph-edit-distance convergence as false at every iteration.
The criterion that actually fired is therefore objective-based, which
weakens the guarantee rather than strengthening it: a structural
oscillation that left $J$ stationary would satisfy it.

\begin{figure}[htbp]
\centering
\includegraphics[width=0.9\textwidth]{leaf_size_hist}
\caption{Distribution of leaf query counts in the frozen snapshot
(construction assignment, not the arena's routing), with the birth
floor $n_{\min} = 55$ marked. The floor is enforced at split time on
the routed partition and is not an invariant of the final tree; the
one leaf of 87 that sits below it (54 queries) occupies the first,
below-floor bin. Median 88, IQR $[71, 118]$, maximum 396.}
\label{fig:leaf-size-hist}
\end{figure}

\paragraph{The evidence per requirement, in one table.}
Table~\ref{tab:requirements-evidence} gives one line of evidence for
each of the four requirements of Section~\ref{sec:arch-requirements},
each with the caveat that bounds it.
One marker used in the table and twice more in this thesis is defined
here once: \textsc{void as to values} means the underlying numbers
were computed on a tree superseded by the frozen artifact and are
withdrawn, while the design of the analysis and any comparison
internal to it survive.

\begin{table}[htbp]
\centering
\small
\caption{Evidence for the four partition requirements of
Section~\ref{sec:arch-requirements}. Every row is ;
the caveat column states what the row does not establish. Full
derivations are in Appendix~\ref{app:dag-formalism}.}
\label{tab:requirements-evidence}
\begin{tabular}{p{0.7cm}p{5.6cm}p{5.8cm}}
\toprule
Req & Evidence & Caveat \\
\midrule
R1
& A split is accepted only when the minimum chance-corrected separation
over all child pairs reaches $0.025$. That bar is positioned between two
measured nulls: an isotropic null whose $p_{95}$ spans $0.0055$--$0.0093$
over $n = 75\ldots900$, and a within-node null whose median lies well
above it.
& The bar is conservative against ``more than nothing'' and permissive
against ``more than this node's own elongation''. The within-node null
values are \textsc{void as to values}; the method survives. \\
\addlinespace
R1
& Against a no-taxonomy floor on the same held-out queries and the same
geometry: the induced 87 leaves score $J = 0.2538$, the editorial
14 MMLU-Pro categories $0.1270$, and those same categories randomly
sub-split to 88 cells $0.1260$. Cell count alone buys nothing.
& A between-partition comparison on held-out top-1 assignment, not a
reproduction of the construction-pool objective. It shows the induced
cells are more separated, not that they are better to judge in. \\
\addlinespace
R1
& Independent of the geometry: rubrics induced on these cells are
specific to their cells against a size-matched randomised null,
Mann--Whitney $U$, one-sided, $p = 1.21\times10^{-6}$
(Section~\ref{sec:res-rubric-null}).
& The measure is lexical. It shows reproduction of a cell's vocabulary,
not application of its criteria. \\
\addlinespace
R2
& 87 leaves; median 88 queries per leaf, IQR $[71, 118]$, minimum 54,
maximum 396.
& \texttt{minClusterSize} is a \emph{birth} constraint enforced at split
time, not an invariant: populations drift under later re-routing and
1 of 87 leaves ends below it. \\
\addlinespace
R2
& Identification is reached at 7 model pairs of 28 --- a spanning tree,
as the theory requires --- within 7 to 10 comparisons per cell in all
52 settled cells, median 7 (Section~\ref{sec:res-cost}).
& Identification only. Resolution needs $\sim$143 comparisons per
pair against the 1.6--5.1 achieved; see the same section. \\
\addlinespace
R2
& Discriminative power is flat across an eleven-fold change in cell
count, so the number of cells can be chosen on budget grounds
(Section~\ref{sec:res-granularity}).
& This is what licenses \texttt{minClusterSize} $= 55$, and it is a
statement about this corpus and this roster, not about partitioning in
general. \\
\addlinespace
R3
& Certified fixed point at iteration 10: edit-phase objective change $0$, maximum
$1 - \cos(\hat{\boldsymbol{\mu}})$ over all nodes
$1.11\times10^{-16}$, \texttt{CERTIFIED true}. Iterations 8, 9 and 10
are identical on every recorded column.
& It certifies \emph{this} artifact, not that the loop terminates from
an arbitrary start --- it does not. The certificate tests period~1
only, and the stop itself came from the objective arm of the
termination rule rather than the structural one. The two objective
fields accumulate in parallel over a float sum and are reproducible
to a bound rather than exactly; no repeat run at this configuration
has been made to measure that bound. \\
\addlinespace
R3
& Across seeds rather than within one: the twenty sweep runs agree on
held-out leaf assignment at an adjusted Rand index of median $0.677$,
range $0.573$--$0.757$ over 190 pairs, against a chance level of $0$.
& Agreement, not equivalence. A fifth to a third of the pairwise
structure differs between any two seeds, and which leaf a given query
lands in is not a seed-invariant fact. \\
\addlinespace
R4
& The acceptance gate is expressed in units of its own objective's
paired standard error. Across the 69 accepted edits that move queries,
$\mathrm{SE}(\Delta J)$ spans $13.4\times$ ($3.27\times10^{-5}$ to
$4.38\times10^{-4}$); the gate declined eight split sites whose
$\Delta J$ cleared $\tau$, four at $z < 1$ and the worst at $z = 0.089$.
& The gate is defended on coherence alone. Both of its empirical
corroborations were tested and refuted (see below). \\
\bottomrule
\end{tabular}
\end{table}

Three of these rows deserve a sentence rather than a table cell.

\textbf{On R1, and on the floor the construction is measured against.}
The objective $J$ has been reported so far only for the artifact, which
says nothing about how much of it the construction earned. Taking the
held-out queries the analysis recorded --- $3{,}575$, a count from an
earlier reserved-pool generation; the settled pool is $3{,}445$, and
the comparison is internal to one fixed set, so the ordering of the
rows does not depend on the generation --- embedding them at the same
$256$-dimensional slice, and scoring four partitions of that same set
with the same estimator gives Table~\ref{tab:j-floor}.

\begin{table}[htbp]
\centering
\small
\caption{The construction objective against a no-taxonomy floor.
All four partitions are scored on one fixed held-out set ($3{,}575$
queries, an earlier pool generation; see the text) at $d = 256$; the two randomised rows are means over 200 draws with
2.5th--97.5th percentile intervals.}
\label{tab:j-floor}
\begin{tabular}{lrr}
\toprule
Partition & Cells & $J$ \\
\midrule
Induced leaves (frozen tree) & 87 & $0.2538$ \\
Induced depth-1 anchors & 14 & $0.1555$ \\
Editorial MMLU-Pro categories & 14 & $0.1270$ \\
Editorial categories, randomly sub-split & 88 & $0.1260$ \ \small $[0.1253, 0.1266]$ \\
Random relabelling at the induced leaf sizes & 87 & $-0.0001$ \ \small $[-0.0026, +0.0022]$ \\
\bottomrule
\end{tabular}
\end{table}

The estimator behaves: a random relabelling into cells of the induced
sizes scores $-0.0001$ with a $95\%$ interval of $[-0.0026, +0.0022]$,
so the chance correction is doing what it claims, and the
implementation used here agrees with the one in the pipeline. The
headline contrast --- separation doubled against the editorial
partition, and none of the gain from cell count alone --- is carried
by the table rows. What the rows do not carry is that the gain
decomposes unevenly: holding granularity at fourteen cells and
replacing the editorial assignment with the tree's own depth-1
anchors is worth $+0.029$; holding the assignment induced and going
from fourteen anchors to eighty-seven leaves is worth $+0.098$. Most
of what the construction earns is earned below the anchor level.

Two limits bound this table. It scores held-out top-1 assignments, so
it is not a re-derivation of the construction-pool objective ---
although the induced row's $0.2538$ against that objective's $0.2531$
is itself evidence that $J$ transfers to unseen queries rather than
fitting the pool it was optimised on. And separation is not
judgeability: the table establishes that the induced cells are more
internally alike than the editorial ones, not that a judge does better
inside them; that question is Section~\ref{sec:res-granularity}.

\textbf{On R4.} The weakest proposal the gate declined carried
$\Delta J = 9.46\times10^{-6}$ against
$\mathrm{SE}(\Delta J) = 1.06\times10^{-4}$; a rule thresholding on
$\tau$ alone would have committed it, along with seven others, and
committing an edit whose measured improvement is a fraction of its own
measurement error is incoherent regardless of what it does to any
downstream number. A second measurement makes the point harder: in the
frozen run the $\tau$ floor never bound at all --- the smallest
threshold any of the $173$ declined proposals was tested against is
$88\tau$ --- so every acceptance decision in the artifact was taken
against a threshold measured from the edit rather than set in advance.
That is the entire argument for the gate, and it is enough: the two
things one would naturally cite in its support (that it tightens
run-to-run stability, and that it raises $J$) were both measured and
both refuted, as the next paragraph reports. The transferable
statement is that \emph{a gate can be right because its units are
right, without improving the number it gates}.

\textbf{On R2.} The cell count is a budget choice, not a statistical
one: \texttt{minClusterSize} $= 55$ was chosen so that the resulting
87 cells match what the judge-call budget can identify, and per-cell
reliability is a consequence of that choice, not its target. The
ordering is defensible only because of the flatness result of
Section~\ref{sec:res-granularity}; without it, 55 would be a
compromise between coherence and support with no evidence behind the
trade.

\paragraph{The acceptance gate is stability-neutral and $J$-neutral.}
both results are negative. The gate's effect on
run-to-run structure is an additive offset of about $4.4$ leaves with
no detectable change in dispersion. The dispersion reading needs its
mechanics stated, because the raw numbers point the wrong way: the
coefficient of variation of leaf counts rose from $14.3\%$ to $15.7\%$
under the gate and was initially read as a degradation, but the
standard deviation is unchanged (variance ratio $1.05$ against an
$F(4,4)$ critical value of $6.39$; the correlation between a seed's
leaf-count reduction and its deviation from the mean is $+0.006$), and
the coefficient of variation rises mechanically because the mean fell.
The correct statement is that no dispersion change is detectable in
either direction, which is the second instance of rule~D3
(Section~\ref{sec:meth-discipline}). And the gate does not raise the
objective it gates: $J$ is higher on 2 of 5 seeds, exact two-sided sign
test $p = 1.000$, median $\Delta J = -0.0001$; the apparent gain was
seed 42 alone. An earlier draft quoted $p = 0.812$ for this
comparison, which no tail of a sign test on five seeds can return; it
is withdrawn.

The gate is adopted on coherence alone (the On~R4
paragraph above), and adopted \emph{after} both of its natural
empirical supports were tested and refused. A corroboration that was
tested and failed is stronger evidence about a rule's justification
than one that was never tested at all.

Caveat: the five-seed sweep's leaf counts predate a correction to the
splitter's routing check and are \textsc{void as to values} (marker
defined above). The analysis design and its
two negative results are unaffected (they are comparisons internal
to the sweep), but the underlying counts need re-deriving.

\subsection{The Adapted Partition, and How Far It Moves}\label{sec:res-adapted-partition}

Chapter~\ref{ch:methodology} defines two partitions of the held-out
set and notes that they coincide only absent cross-domain migration.
This subsection measures the migration, because every per-domain
number in this chapter is stated over the adapted partition and the
two do not agree.

Table~\ref{tab:anchor-cells} gives, for each of the
fourteen anchors, the structural cell count $|\mathcal{L}(D_k)|$, the
held-out pool assigned to it, the median cell size, and the fraction of
that pool whose MMLU-Pro label matches the anchor. The cell counts sum
to 87 and the pools to 3{,}445, so the table is a partition of the
frozen artifact and of the held-out set, with no double-counting and no
remainder. For the eight domains that ran as arenas, the cell count in
this table equals the number of cells the run actually scheduled
comparisons in, checked run by run.

\begin{table}[htbp]
\centering
\small
\caption{The adapted partition of the frozen artifact. Cells are
$\mathcal{L}(D_k)$, the leaves beneath anchor $D_k$; the pool is the
held-out questions whose primary leaf lies in that subtree, which is
the scoping rule every paired run applies. Label agreement is the
fraction of the pool carrying the matching MMLU-Pro category, and so
measures how far the induced partition reproduces the editorial one.
Rows marked $\ast$ ran in the settled batch, and \emph{run pool} is
the in-scope question count that run recorded at routing time. It is
a larger quantity than the \emph{judged} question count of
Table~\ref{tab:run-record}, because a run's budget can end before
every in-scope question receives a comparison (engineering routed
$202$ and judged $158$). Cells, pool, median and label agreement come
from the routing-harness export
(\texttt{reserved\_leaf\_assignments.csv}); the run pool comes from the
runs themselves, and the two routing paths do not agree --- see the
note below the table.}
\label{tab:anchor-cells}
\begin{tabular}{lrrrrr}
\toprule
Anchor & cells & pool & median cell & label agreement & run pool \\
\midrule
Mathematics$^\ast$      & 11 & 334 & 27 & 81.4\% & 291 \\
Physics$^\ast$          & 10 & 338 & 32 & 75.1\% & 237 \\
Chemistry               &  8 & 401 & 38 & 72.6\% & --- \\
Psychology$^\ast$       &  8 & 258 & 28 & 67.8\% & 265 \\
Other                   &  7 & 205 & 28 & 60.0\% & --- \\
Biology                 &  6 & 220 & 29 & 70.9\% & --- \\
Computer science$^\ast$ &  6 & 191 & 31 & 49.7\% & 159 \\
Economics               &  6 & 236 & 34 & 83.1\% & --- \\
Law$^\ast$              &  6 & 280 & 46 & 94.6\% & 277 \\
Health                  &  5 & 191 & 35 & 71.7\% & --- \\
Engineering$^\ast$      &  4 & 250 & 42 & 72.8\% & 202 \\
History$^\ast$          &  4 & 133 & 38 & 71.4\% & 131 \\
Business                &  3 & 241 & 90 & 71.4\% & --- \\
Philosophy$^\ast$       &  3 & 167 & 46 & 72.5\% & 130 \\
\midrule
Total                   & 87 & 3{,}445 & --- & --- & --- \\
\bottomrule
\end{tabular}
\end{table}

\paragraph{Two routing paths, and where they disagree.} The cell counts
are a property of the tree and are not in doubt: they are read from the
frozen snapshot and independently confirmed by every run. The
\emph{pools} are not so clean. The harness export routes cached
embeddings through \texttt{TaxonomyTrickler} and assigns all 3{,}445
held-out questions; each arena run re-routes live and discards 561 of
them as outliers matching no leaf, the same 561 in all eight runs. The
remainder does not merely shrink each pool proportionally --- physics
falls from 338 to 237 while psychology rises from 258 to 265 --- so the
two paths disagree about composition, not only about coverage. The
disagreement was found while reconciling the leaf-count definitions and
is \textsc{not diagnosed}. It bounds nothing in the paired contrasts,
which compare two arms over one identical question set drawn by one
routing path, but it does mean the pool column should be read as the
frozen artifact's assignment and not as the set any run judged. Each
run's own count is the last column, and each run reports its own
sampled coverage in the subsection that presents it.

Label agreement runs from $49.7\%$ to $94.6\%$.
The induced partition is therefore not a relabelling of the editorial
one, which is the intended behaviour --- construction sees embeddings
and never sees the category field --- but it has two consequences that
bear on how the arena results are read, and both cut against the
partition claim.

First, a domain's arena run is only partly about that domain as
MMLU-Pro defines it. Law is nearly pure at $94.6\%$ and mathematics is
high at $81.4\%$, so their runs are close to what their names suggest.
Computer science is the extreme case: barely half its pool carries the
computer-science label, and its six cells include leaves whose
subject matter is algebraic structure, statistical inference and time
series. A result on that pool is a result about a region of the
embedding space that the construction anchored to computer science,
not about computer science as a syllabus topic, and it is reported on
that basis.

Second, the domain screen of Section~\ref{sec:res-law} is computed
over the \emph{canonical} partition, because it uses ground-truth
accuracy by MMLU-Pro category and needs no taxonomy at all. The arena
runs it is compared against are scoped by the adapted partition. The
two are being asked about overlapping but distinct question sets, and
the mismatch is largest exactly where label agreement is lowest. This
is one more reason, independent of the ones in
Section~\ref{sec:res-law}, not to treat the screen as a clean
predictor of which domains reward a partition.

\paragraph{A correction to Addendum 3.} The registration for physics
and computer science described them as having 12 and 3 leaves, with
pools of 546 and 208 questions. Those figures come from the
label-composition count that Chapter~\ref{ch:methodology} sets aside:
they assign each leaf to its majority MMLU-Pro category rather than to
its anchor. Under the adapted partition the runs actually used, the
figures are 10 and 6 cells over pools of 338 and 191. The error is
descriptive and touches no prediction --- both registrations turn on
the sign and magnitude of $\Delta\rho$, neither of which references a
cell count. The secondary observation registered for physics also
survives: mathematics still has the most cells, physics still has the
second most, and the gap is one cell under either count. The
correction is recorded here rather than repaired in the registration,
which is left as written.

\subsection{Link 2: Rubric Specificity Against a Randomised Null}\label{sec:res-rubric-null}

The question is whether coherent cells produce more specific rubrics
than incoherent ones. It is offline: it needs no arena. It was
pre-registered with four directional predictions, the confound
direction named per measure, and decision bands fixed in advance.

Rubrics were induced by the identical procedure on
the 87 leaves and on 13 size-matched random cells drawn from the same
corpus. The specificity measure (the fraction of a rubric's
vocabulary found in its own cell's queries, above what it shares with
other cells) separates the two arms decisively: Mann--Whitney $U$,
one-sided, $p = 1.21 \times 10^{-6}$, with medians $0.138$ against
$0.012$. The result holds under all three definitions of the
per-cell statistic that were considered
($1.2\times10^{-6}$, $2.1\times10^{-6}$, $7.6\times10^{-6}$).

All four pre-registered discriminators fire in the
predicted directions (Table~\ref{tab:rubric-null}).
Figure~\ref{fig:rubric-specific} puts that separation next to its
consequence at the verdict.

\begin{figure}[htbp]
\centering
\includegraphics[width=\textwidth]{rubric_specificity_null}
\caption{The rubrics are specific to their cells: cumulative
distributions of own-cell rubric specificity for the 87 induced cells
against 13 random cells drawn to the same sizes and put through the
same induction procedure, medians $0.138$ against $0.012$,
Mann--Whitney $U$, one-sided, $p = 1.21 \times 10^{-6}$. Zero on the
axis means a rubric says no more about its own cell than about any
other. The analysis is offline --- rubric text against rubric text ---
and involves no arena verdict; that the specificity is nonetheless
inert at the verdict is Section~\ref{sec:res-c5}.}
\label{fig:rubric-specific}
\end{figure}

\begin{table}[htbp]
\centering
\caption{Rubric-specificity discriminators, leaf cells against
size-matched random cells. The predicted direction was fixed in advance
for each row.}
\label{tab:rubric-null}
\begin{tabular}{lrrl}
\toprule
Measure & 87 leaf cells & 13 random cells & predicted \\
\midrule
own-cell specificity (median) & \textbf{0.138} & 0.012 & leaf higher \\
pairwise Jaccard between rubrics & 0.0754 & 0.1187 & leaf lower \\
terms present in $\geq 90\%$ of rubrics & \textbf{0} & 6 & leaf $\approx 0$, random $>0$ \\
cells with positive specificity & 81/87 (93\%) & 7/13 & leaf higher \\
\bottomrule
\end{tabular}
\end{table}

The sharpest of these is the third. A random cell contains chemistry,
law and history at once, so its rubric must fall back on domain-general
language; a generic core is present across the random rubrics and
absent from every leaf rubric. Had both been near zero the instrument
could not see the difference and the test would have been inconclusive
rather than negative, a distinction fixed in advance.

The random arm's absolute specificity of $0.012$ falls inside the
$\leq 0.05$ band fixed before the run as meaning ``coherence is the
cause''. An arm-size control was run because 13 against 87 is
unbalanced: 13 leaf rubrics subsampled give $0.134$, and a 200-draw
bootstrap over subsamples gives $0.138$, so the treatment arm's value
is not a function of its size.

One consequence belongs here rather than in
Chapter~\ref{ch:discussion}. These are the same rubrics whose
replacement by a generic prompt leaves the judge's ranking unchanged
in every settled domain (Section~\ref{sec:res-c5}). Link~2 holds, the
cells do produce measurably specific rubrics, and on this corpus that
specificity is inert at the verdict.

Four qualifications travel with the result and none is optional.
\begin{enumerate}
\item \textbf{The measure is lexical throughout.} It demonstrates
reproduction of a cell's vocabulary, not application of its criteria.
\item \textbf{The treatment bundles two things.} Coherent membership
comes with an informative topical label, and the null cells were
deliberately labelled \texttt{Cluster NN}, so the design supports ``a
coherent cell together with its label causes specificity'' rather than
coherence alone.
\item \textbf{The null arm is thin, and thinner than registered.} It
was pre-registered at 87 random cells plus a 14-cell domain-level
intermediate arm, and run at 13 size-matched random cells with the
intermediate arm not run at all. The deviation is a budget decision
taken after registration, and it is recorded here because the
separation would read as stronger if the registered design were
assumed.
\item \textbf{The supporting artifacts are untracked.} The null-arm
generation script and its outputs sit in an untracked build directory.
The headline separation was reproduced from those stored artifacts on
2026-07-27 and re-derived independently on 2026-07-30, robust across
all three specificity definitions (every $p < 10^{-5}$). A dated
addendum (2026-07-28) records both the written-back results and this
caveat, because a result whose artifacts are untracked is one clean
checkout from being unsupported.
\end{enumerate}

The same addendum resolves which decision band governs which quantity,
because the two pre-registration documents state bands for different
things. The $\geq 0.05$-with-$\geq 75\%$-of-cells-positive band governs
the treatment arm's within-arm specificity, met at $0.138$ with $93\%$
of cells positive. The $\leq 0.05$ / $\geq 0.10$ band governs the
random arm's absolute own-cell specificity, met at $0.012$, and is the
band quoted above.

\subsection{Link 3's Premise Fails, and It Was Pre-Registered}\label{sec:res-granularity}

Cells do not carry different true model rankings on
this corpus, at any granularity from 14 to 152 cells. The analysis was
fixed in writing before it was run (three outcomes named in advance,
``observed $\rho$ is primary''), and that
registration is the only reason the result is readable, because the
disattenuated variant of the statistic (reported in the analysis
export, not tabulated here) moves the opposite way and clips past
$1.0$.

The question is whether cells carry different \emph{true} model
rankings: for every pair of cells, the Spearman correlation between the
two cells' ground-truth accuracy rankings of the roster. If cells
differ in what they measure, that correlation should fall as cells get
finer. Table~\ref{tab:granularity} reports the medians at four
granularities.

\begin{table}[htbp]
\centering
\caption{Observed between-cell Spearman $\rho$ of ground-truth model
rankings, across an eleven-fold change in granularity, under two
assignment conventions (A) and (B) defined in the note below.
Ground truth only --- no judge verdict enters this table, so the run
withdrawals do not touch it; it was computed at the pre-settlement
rosters (the historical eight-model band and the eleven-model band),
and the settled-roster leaf-substructure refit of
Section~\ref{sec:res-leaf-null} confirms the flatness it shows. The
88-leaf and 87-leaf rows differ only in the router used to assign
queries and are an explicit control.}
\label{tab:granularity}
\begin{tabular}{lrrrr}
\toprule
Partition & cells & cell pairs & median $\rho$ (A) & median $\rho$ (B) \\
\midrule
MMLU-Pro editorial domains & 14 & 91 & 0.929 & 0.929 \\
induced leaves, pre-correction snapshot & 88 & 3828 & 0.922 & --- \\
induced leaves, canonical snapshot & 87 & 3741 & 0.922 & 0.884 \\
induced leaves, finer setting & 152 & 11476 & 0.905 & --- \\
\bottomrule
\end{tabular}
\par\smallskip
\footnotesize Two assignment conventions, because the choice changes
the leaf-level value and not the conclusion. Column (A) uses the
snapshot's own construction assignments over the full corpus, so the
$0.922$ in it is \emph{not} a re-routed value; column
(B) re-routes the reserved pool through the production router and
correlates only over held-out queries
(\texttt{tools/analysis/discriminative\_flatness.py},
\texttt{reserved\_leaf\_assignments.csv}, 87 leaves and 3{,}863
assignments). Both are defensible readings of ``the ranking within a
leaf'' and they are not the same quantity. Flatness holds under either,
since each column compares rows computed the same way; the 88-leaf and
152-leaf rows were produced under convention (A) only and were not
regenerated under (B).
\end{table}

It does not fall. Interquartile ranges overlap throughout. The 87- and
88-leaf rows are identical to three decimals although they differ in
the router used to assign queries, so discriminative power is flat
across \emph{routers} as well as across granularities.

Centring out each model's global mean is a control on that flatness,
and it is the weakest measurement in this section. If cell identity
carried nothing, the centred between-cell correlations should sit at
the mechanical centring null $-1/(C-1)$, which is $-0.077$ at 14
domains and $-0.011$ at 87 leaves. What the script\footnote{%
\texttt{tools/analysis/discriminative\_flatness.py}.} returns is
$+0.024$ at 14 domains and $-0.024$ at 87 leaves on the 8-model roster,
and $-0.109$ and $-0.009$ on the 11-model band. The residual is small
everywhere, and its sign is unstable across rosters at the 14-domain
cut. At that cut the 8-model value is nearer zero than the null is, so
this control cannot separate ``zero'' from ``the centring null'' in
either direction, and it is reported as small rather than as located.
An earlier draft reported a single quartet $-0.119$, $-0.024$,
$-0.024$, $-0.024$; no one centring convention reproduces all four
values, so the quartet is withdrawn and the four roster-specific values
above are reported in its place. They have not themselves been
independently confirmed.

The join underlying this analysis was verified before it was computed,
per rule~D2: exact, $11769/11769$; all roster models present in every
cell of every tree; no cells dropped.

At the corrected roster. The flatness survives the
change of roster. Re-derived on re-routed, reserved-only assignments
by the same script, the 14-domain and
87-leaf medians are $0.929$ and $0.884$ at the 8-model pilot roster
against $0.948$ and $0.908$ at the 11-model length-matched band. The
granularity drop is $0.045$ at one roster and $0.040$ at the other, and
the interquartile ranges overlap in both cases ($[0.857, 0.970]$
against $[0.806, 0.934]$; $[0.927, 0.964]$ against $[0.860, 0.942]$).
An earlier draft reported $0.973$ and $0.936$ for the 11-model band.
Those two values could not be regenerated by any script or
configuration in this repository and are withdrawn under the evidence
convention of Section~\ref{sec:research-questions}; the substitution
does not change the finding, since both rosters give the same small
drop with overlapping spread.

\paragraph{Three bounds, all pre-registered.}
\emph{Power, not proof.} Per-model per-cell accuracy has a standard
error of about $2.1$ percentage points at 14 domains, $5.8$ at 87
leaves and $6.9$ at 152. The null is consistent with ``no cell-specific
ranking signal'' and with ``a signal below about 6 points at leaf
granularity''; this analysis cannot separate them. The 14-domain row is
the stronger evidence precisely because it is the low-noise condition,
and its centred residual is small there too.
\emph{Roster.} The rosters this table was computed at span $54$ to
$71$ accuracy points depending on the pool, and the settled roster
spans $74$ (Table~\ref{tab:roster-r8}). At any of those spreads,
$\rho \approx 0.92$ was largely measuring that the strongest model beats
the weakest in every cell. A roster clustered within about 10 accuracy
points is untested and is the version of this question that could still
come out positive.
\emph{Ceiling, not prediction.} This is measured on each cell's own
constituent queries and on ground-truth accuracy, never on routed
held-out queries and never on judge verdicts. It \emph{bounds} what an
arena could discover in these cells; an arena can show less
discrimination than this, never more.

This is a property of the corpus and the roster, not
of the method. MMLU-Pro is saturated for a roster of this spread, and
cell identity contributes essentially nothing to model ranking beyond
global capability ordering, at any granularity from 14 to 152, under
either router. It says nothing about judge quality, which is what
Section~\ref{sec:res-judge} measures.

No figure accompanies Table~\ref{tab:granularity}, and the omission is
deliberate. A box plot of four heavily overlapping distributions makes
a null look like a finding, because the eye reads the small median
offsets and not the overlap that carries the result. The plotting
script written for it declines to draw.

\subsection{Cells Do Not Create Evaluative Diversity}\label{sec:res-leaf-null}

The granularity analysis asks whether between-cell agreement falls as
cells get finer. A sharper question is whether any individual domain's
leaves disagree with each other by more than sampling noise, and it
needs its own null, because leaf sizes here are 21 to 88 questions. A
model's accuracy on one leaf carries a standard error near 8 to 9
points while adjacent models in this roster differ by 2 to 4, so two
noisy rankings of near-tied models can disagree strongly with nothing
underneath.

Ground truth only, no judge and no arena, at the
settled eight-model roster. For each domain, that domain's reserved
questions are reshuffled at random into pseudo-leaves of \emph{exactly}
the observed leaf sizes, and the same between-leaf Spearman statistic
is recomputed; 2{,}000 draws, seed 42.\footnote{Conventions of
\texttt{tools/analysis/leaf\_substructure\_null.py}.} Real leaves that
carry sub-structure should sit below the shuffled distribution, at
lower agreement than chance. \textbf{No domain does}: the smallest of
the fourteen empirical $p$-values is $0.122$, against a Bonferroni
threshold of $0.0036$. Nor does any individual leaf pair: of the $293$
within-domain pairs, \textbf{zero} survive correction at this roster.

Both findings were sharper at the superseded twelve-model roster ---
three mathematics leaf pairs, all separating continuous from discrete
mathematics, then survived correction --- and the difference is the
roster, not the corpus. A spread of adjacent-model gaps wide enough to
make pairs key-decidable (Section~\ref{sec:res-correctness}) is also
wide enough that no leaf-sized sample reorders it; whatever fine
sub-structure the corpus carries sits below what this roster can
express. The claim this section supports is therefore bounded: at the
roster the arena actually ranks, cells create no measurable evaluative
diversity.

It is the fifth error this apparatus caught in
its own numbers. The first version of this null sampled pseudo-leaves
\emph{independently}, which gave them 11--26\% pairwise question
overlap where real leaves have 1.6--4.9\%. Overlapping pseudo-leaves
agree with each other for free, so the null was inflated and the real
cells looked distinctive against it. The corrected version matches
pseudo-leaf overlap to the real distribution and returns the opposite
answer. The first version's conclusion is withdrawn
(Appendix~\ref{app:corrections-register}); this is rule~D1 of
Section~\ref{sec:meth-discipline} applied to a comparison arm that had
never itself been checked.

It locates where the boundary of that claim
runs. In the settled roster no analysis of near-tied pairs is even
possible: the smallest aggregate accuracy gap between any two of its
models is $7.9$ points, so there are no ties for a domain to reorder.
The wider corpus does contain the phenomenon --- among the $36$ valid
models, pairs within $3$ aggregate points exist that reverse by $15$
to $22$ points on individual MMLU-Pro
domains\footnote{\texttt{tools/analysis/export\_gt\_scores.py}.} ---
so domain-level
parity can hide margin-level disagreement in a clustered roster. On
\emph{this} roster it cannot, by construction.

Together with the granularity analysis, this is what
forces any apparent per-cell variation in the batch to be read as
sampling noise rather than as discovered specialisation --- and it is
also the honest boundary of the negative result. Cells create no
evaluative diversity \emph{for a roster whose models are far apart};
whether they would for a tightly clustered roster is untested and is
named as future work (Section~\ref{sec:disc-future}).

\subsection{The Domain Screen, on All Fourteen Domains}\label{sec:res-law}

An arena can only reveal cell-specific capability where cell-specific
capability exists. The screen is offline: for each MMLU-Pro domain,
correlate the roster's ground-truth accuracy ranking within that domain
against its global ranking, and compare against a size-matched
permutation null over random subsets of the same reserved pool. It
costs no judge calls. It was first computed, on all fourteen domains
and before any arena was scheduled, at the since-superseded
twelve-model roster; that fit is what fixed the eight-domain selection,
and it flagged four domains as reordering. The historical fit and the
selection it produced are recorded in the pre-registration
(Appendix~\ref{app:preregs}); its $p$-values are not carried forward.

Refit at the settled eight-model roster: \textbf{the
screen saturates} (Table~\ref{tab:domain-screen}). The domain-level
ground-truth ordering equals the global ordering exactly
($\rho = 1.000$) in twelve of fourteen domains. Only two depart at all:
law ($\rho = 0.994$, $p_{\text{emp}} = 0.003$) and chemistry
($\rho = 0.976$, $p_{\text{emp}} < 0.001$), each by a single adjacent
transposition of nearly-tied models. The roster property that makes
model pairs key-decidable --- an accuracy spread of $16.5\%$ to
$90.5\%$ --- is the same property that makes domain membership nearly
irrelevant to the ordering: models this far apart globally are this far
apart everywhere.

\begin{table}[htbp]
\centering
\caption{The offline reordering screen on all fourteen MMLU-Pro
domains, refit at the settled eight-model roster against a size-matched
1{,}500-draw permutation null (seed 42, reserved split, 3{,}437 covered
questions). A \emph{low} $p$ means the domain's ground-truth ranking of
the roster departs from the global ranking more than a random subset of
the same size does. Domains marked $\bullet$ are the eight the arena
ran; the selection was fixed under the earlier twelve-model fit and
carries no endorsement from this one.}
\label{tab:domain-screen}
\begin{tabular}{lrrl}
\toprule
Domain & $\rho$ & $p_{\text{emp}}$ & run \\
\midrule
\multicolumn{4}{l}{\textit{Reorders, marginally}} \\
chemistry        & 0.976 & $<0.001$ & \\
law              & 0.994 & 0.003 & $\bullet$ \\
\midrule
\multicolumn{4}{l}{\textit{Saturated: domain ordering equals the global ordering}} \\
biology          & 1.000 & --- & \\
business         & 1.000 & --- & \\
computer science & 1.000 & --- & $\bullet$ \\
economics        & 1.000 & --- & \\
engineering      & 1.000 & --- & $\bullet$ \\
health           & 1.000 & --- & \\
history          & 1.000 & --- & $\bullet$ \\
math             & 1.000 & --- & $\bullet$ \\
other            & 1.000 & --- & \\
philosophy       & 1.000 & --- & $\bullet$ \\
physics          & 1.000 & --- & $\bullet$ \\
psychology       & 1.000 & --- & $\bullet$ \\
\bottomrule
\end{tabular}
\par\smallskip
\footnotesize Conventions of
\texttt{tools/analysis/domain\_reorder\_screen.py} (Spearman with tie
ranks; one-sided empirical $p$, fraction of null draws at or below the
observed $\rho$). At $\rho = 1.000$ the empirical $p$ is uninformative
--- the majority of null draws also reach $1.000$ --- and is not
reported.
\end{table}

The consequence for reading the batch is stated once, here: the eight
domains function as \textbf{replications of one measurement}, not as
screened contrasts. This is the same fact as the shared ground-truth
ordering of Section~\ref{sec:res-correctness}, seen from the screen's
side, and it is why no per-domain arena result in this thesis is read
as evidence about that domain in particular.

The reason this thesis computes the screen at
the roster the arena actually runs. \textbf{At earlier rosters the
screen reversed twice, and both reversals had the same cause.} At an
early 8-model band, math was indistinguishable from a same-size random
subset and the recommendation was law. After the trace backfill the
screen was re-run at 37 models, math became significant and the
recommendation flipped to math --- a flip resting on a model
intersection poisoned by a single model in a different
question-identifier space, whose inclusion collapsed law from 287
questions to 20. Excluding it restored law. Both inversions came from
evaluating a ranking statistic at a model count that did not match the
decision it was informing; the saturation of the settled refit is the
same lesson from the other side, a screen answering only questions
about the roster it is given.

Two caveats bound the whole section. First, the screen uses MMLU-Pro's
native category labels rather than induced cells, so it bounds what a
partition of a domain could show without showing that induced cells
inside any domain separate. Second, and measured at the earlier
roster, it does not even attempt the within-domain question: the
correlation between a domain's screen $p$ and the agreement between
its own leaves was $+0.072$ (Section~\ref{sec:res-leaf-null}). The
screen is a filter on a domain aggregate, and no claim in this thesis
rests on it as more than that.

\subsection{Polyhierarchy: A Measured Negative Result}\label{sec:poly-negative}

The construction produces a strict tree: every node below the root has
exactly one parent. An earlier version of the pipeline admitted
polyhierarchy (a \emph{concept} belonging to several branches)
through a dedicated cross-link operator that gave an existing node a
second parent when the new parent's residual pool demonstrably fitted
it. The operator is gone from the construction described in this
thesis; the single-parent invariant that replaced it, and the one
restriction on fusion needed to hold it, are specified in
Appendix~\ref{app:cross-link-operator}. Construction-time routing
remains weighted multi-membership (Appendix~\ref{app:dag-formalism});
that is a claim about \emph{queries}, not concepts, and it is
retained. What follows is the evidence that retired the operator.

Across 42 cross-link proposals on the canonical
corpus, the capture fraction $f$ that the operator uses to rank
candidates was uncorrelated with the objective it must satisfy:
$r(f, \Delta J) = +0.085$, with complete distributional overlap between
accepted and rejected proposals (median $f$ of $0.802$ accepted against
$0.818$ rejected).
The effect sizes were an order of magnitude apart: median $\Delta J$
was $6.9 \times 10^{-5}$ for cross-links against $5.3 \times 10^{-4}$
for splits.
Admission was therefore decided by evaluation order rather than by
evidence, which the proposal sequence makes explicit: one node was
accepted at four hosts and rejected at two, and acquired five parents
by iteration~1.

Cost of removal: the multi-leaf rate of
held-out queries was $0.1549$ with cross-links and $0.1518$ without,
so queries that legitimately span domains still reach several leaves
through the routing beam and the membership floor.

$J$ scores a partition of query mass into cells.
Giving a node a second parent changes which paths reach its cell, but
it leaves the cell contents almost unchanged, so $\Delta J$ carries
little signal about whether the edge is semantically right.
A partition-based objective is structurally indifferent to
polyhierarchy, and no threshold on such an objective can separate good
cross-links from bad ones. The operator was therefore removed rather
than retuned. What is given up is the concept-level claim, which the
evidence above could not support in the first place.

One caveat on provenance, stated because it bounds what the numbers in
this subsection can be checked against. They come from a construction the
current source can no longer produce: the operator was deleted, and with
it the proposal type its records were written under, so the proposal logs
of the frozen artifact and of every run that postdates the deletion
contain no cross-link rows. The measurements stand as a record of why the
operator was retired and cannot be re-derived from any artifact this
thesis reports.

\section{Chapter Summary}\label{sec:res-summary}

Stated against the criteria fixed before the data, with no
interpretation: Table~\ref{tab:res-summary} is the record, and the
sections in its last column carry every qualification.

\begin{table}[htbp]
\centering
\small
\caption{Chapter summary. Verdicts are stated against the criteria
fixed before the data existed.}
\label{tab:res-summary}
\begin{tabular}{p{1.4cm}p{7.8cm}p{3.2cm}}
\toprule
RQ / link & Verdict & Where \\
\midrule
RQ1 & The judge's verdict is produced jointly by answer verification
and a competence prior: the settled batch recovers the key's ordering
on $196/200$ separable pairs across eight domains --- one shared
ordering, recovered eight times --- with $86\%$ agreement on
discriminable comparisons where the judge commits; agreement splits
on candidate strength ($91.1\%$ where the correct model is the
pair's stronger, $56.5\%$ where it is the weaker); the
tie rate nearly doubles ($13.0\%$ to $24.0\%$) where the key stops
discriminating, in all eight domains; and off-subject questions are
judged at on-subject rates. The tie convention moves no $\rho$ in the
batch; the scheduler seed moves $\Delta\rho$ by six grid steps, so
only per-verdict statistics carry decisions. &
\S\ref{sec:res-correctness}--\ref{sec:res-tie-policy}; the split:
\S\ref{sec:res-competence} \\
\addlinespace
RQ2 & The cell rubric changes decisiveness, not ranking. Both arms
produce the same ordering in seven of eight domains (one adjacent
transposition in the eighth); the cell-scoped arm abstains on
$13.0\%$ of decidable comparisons against $10.9\%$, in all eight
domains; accuracy conditional on committing differs by one point
($86.0\%$ against $87.2\%$, both-committed sign test $11/26$,
$p = 0.020$, no single domain significant). Attributable to the arm
bundle, not the rubric text alone. &
\S\ref{sec:res-c5} \\
\addlinespace
Link 1 & Met. Certified fixed point at iteration 10; 87 leaves, median
size 88; the acceptance gate is stability-neutral and $J$-neutral,
both measured. &
\S\ref{sec:res-construction} \\
\addlinespace
Link 2 & Positive. Rubrics induced on coherent cells are specific to
those cells against a size-matched randomised null,
$p = 1.21 \times 10^{-6}$, all four pre-registered discriminators in
their predicted directions. Lexical measure only. &
\S\ref{sec:res-rubric-null} \\
\addlinespace
Link 3 & Premise refused under pre-registration, before the arena ran:
between-cell agreement of true rankings is flat across an elevenfold
change in granularity and across two routers, bounded by power, by
roster spread, and by its ceiling character. At the settled roster the
overlap-matched permutation null is exceptionless --- no domain, and
zero of 293 leaf pairs, survives correction; the three mathematics
pairs that once did belong to the superseded, more clustered roster
and locate the boundary. &
\S\ref{sec:res-granularity}, \ref{sec:res-leaf-null} \\
\bottomrule
\end{tabular}
\end{table}

Polyhierarchy sits outside the three links: operator built, measured,
retired. Query-level multi-membership survives; the concept-level
claim does not (Section~\ref{sec:poly-negative}).
